\documentclass[10pt]{article}
\usepackage[margin=0.78in]{geometry}
\usepackage{amsmath,amssymb,amsthm,mathtools}
\usepackage{graphicx}
\usepackage{xcolor}
\usepackage{microtype}
\usepackage[hidelinks]{hyperref}
\setlength{\parindent}{0pt}
\setlength{\parskip}{5pt}
\newtheorem{theorem}{Theorem}
\newtheorem{proposition}[theorem]{Proposition}
\newtheorem{corollary}[theorem]{Corollary}
\newcommand{\N}{\mathbb N}
\newcommand{\Q}{\mathbb Q}
\newcommand{\F}{\mathcal F}
\newcommand{\W}{\mathcal W}
\newcommand{\ellone}{\ell_1}

\title{Rational Weak-Null Witnesses Generate Every\\
Non-Almost-Schur Filter\\
\large An explicit basis for the negation half of the problem in arXiv:0903.0659}
\author{Automated proof attempt\\
\texttt{agent\_lane\_03} \quad model: GPT5.6}
\date{11 August 2026}

\begin{document}
\maketitle

\begin{center}
\fcolorbox{orange!70!black}{orange!5}{\parbox{0.92\linewidth}{
\textbf{Status: substantial partial result, likely valid, pending human
verification.}  A canonical rational family is an explicit basis for all
filters that are \emph{not} almost Schur, fully answering the second half of
the source request.  The requested basis for the decreasing almost-Schur
class itself remains open.}}
\end{center}

\section*{Source problem and definitions}

All filters below are free filters on $\N$.  A filter $\mathcal G$
\emph{dominates} $\F$ when $\mathcal G\supset\F$.  The source calls $\F$
\emph{almost Schur} if every sequence $(x_n)\subset\ellone$ that converges
weakly to zero along $\F$ has norms not separated from zero.  Thus $\F$ is
not almost Schur exactly when some weakly $\F$-null sequence satisfies
\(\inf_n\|x_n\|_1>0\).

The non-almost-Schur property is increasing under domination.  A subclass is
a \emph{basis} for an increasing class when its upward closure is the whole
class.  Section 6 of the source asks for explicit bases for the decreasing
almost-Schur class and, separately, its increasing negation class
\cite{source}.  The appended page-15 crop gives the exact wording.

\section*{Proof intuition}

Failure of the almost-Schur property comes with a witness sequence.  The
filter contains every index set on which finitely many weak tests of that
sequence are small.  Those weak-neighborhood traces therefore generate a
smaller canonical bad filter.

There seems to be uncountable analytic data in the witness, but its vectors
can be normalized and approximated term by term by finitely supported rational
unit vectors.  The error tends to zero in norm, hence along every free filter.
This turns every failure into a completely coded rational failure without
changing any of its weak-neighborhood traces up to enlargement.  Taking all
such rational trace filters gives exactly the requested upward basis.

\section*{The canonical filters}

Let
\[
 S_{\Q}=\{u\in c_{00}(\Q):\|u\|_1=1\}.
\]
Call a sequence $u=(u_n)_{n\in\N}\subset S_{\Q}$ \emph{admissible} if for
every $k,N\in\N$, every $\phi_1,\dots,\phi_k\in\ell_\infty$, and every
$\varepsilon>0$, the set
\begin{equation}\label{eq:trace}
 V(u;\phi_1,\dots,\phi_k,\varepsilon,N)
 =\{n\ge N:\max_{1\le j\le k}|\phi_j(u_n)|<\varepsilon\}
\end{equation}
is nonempty.  Equivalently, zero belongs to the weak closure of every tail of
$(u_n)$.

For admissible $u$, finite intersections of sets in \eqref{eq:trace} contain
another such set: take the union of the finite test families, the smaller
tolerance, and the larger tail index.  Hence these sets generate a proper
filter, denoted $\W_u$.  It is free because every generator lies beyond an
arbitrary prescribed $N$.  By construction,
\[
 u_n\longrightarrow_{\W_u}0\quad\text{weakly},
 \qquad \|u_n\|_1=1.
\]
Thus $\W_u$ is not almost Schur.

Admissible sequences exist.  Indeed, enumerate the countable dense set
$S_{\Q}$ so that every member occurs infinitely often.  Given finitely many
functionals, their common kernel in $\ellone$ is nonzero; approximate a unit
vector in that kernel by a member of $S_{\Q}$.  A sufficiently accurate
approximation lies in the prescribed finite weak neighborhood and recurs in
every tail.

\section*{Exact basis theorem}

\begin{theorem}[Rational witness basis]\label{thm:basis}
The class
\[
 \mathfrak B_-:=\{\W_u:u\subset S_{\Q}\text{ is admissible}\}
\]
is a basis for the increasing class of all non-almost-Schur filters.
Equivalently, for every free filter $\F$,
\begin{equation}\label{eq:criterion}
 \F\text{ is almost Schur}
 \quad\Longleftrightarrow\quad
 \W_u\not\subset\F\text{ for every admissible }u.
\end{equation}
\end{theorem}

\begin{proof}
We already observed that every $\W_u\in\mathfrak B_-$ is not almost Schur.
Since that property is increasing, every filter dominating a member of
$\mathfrak B_-$ is also not almost Schur.

Conversely, suppose $\F$ is not almost Schur.  Choose a weakly
$\F$-null sequence $(x_n)\subset\ellone$ and $\delta>0$ such that
$\|x_n\|_1\ge\delta$ for every $n$.  Put
\[
 y_n=\frac{x_n}{\|x_n\|_1}.
\]
The scalar multipliers are bounded by $1/\delta$, so $(y_n)$ is still weakly
$\F$-null and every $\|y_n\|_1=1$.

The rational finite-support unit sphere is norm dense in the unit sphere of
$\ellone$.  Choose $u_n\in S_{\Q}$ with
\(\|u_n-y_n\|_1\le1/n\).  The error sequence is norm null in the ordinary
sense and therefore norm null along the free filter $\F$.  It follows that
$(u_n)$ is weakly $\F$-null.

Every set in \eqref{eq:trace} therefore belongs to $\F$: finite weak
convergence gives the same trace without the tail restriction, and freeness
adds $\{n:n\ge N\}$.  In particular, the sequence $u$ is admissible and the
filter generated by all its traces satisfies $\W_u\subset\F$.  Hence every
non-almost-Schur filter dominates a member of $\mathfrak B_-$, proving both
the basis assertion and \eqref{eq:criterion}.
\end{proof}

The construction is fully combinatorial if each $u_n$ is viewed as a finite
rational signed measure on $\N$ and each $\phi\in\ell_\infty$ as a bounded
scalar sequence.  No arbitrary vectors remain in the basis generators.

\section*{Why the positive half does not follow by maximal extension}

Criterion \eqref{eq:criterion} reduces the missing half to finding a useful
explicit cofinal family of filters that simultaneously avoid every
$\W_u$.  Individual blockers exist, but their complements need not retain the
finite-intersection property.  The most obvious Zorn argument also fails
because the almost-Schur class is not closed under unions of chains.

\begin{proposition}[Chain obstruction under CH]\label{prop:chain}
Assume the continuum hypothesis.  There is a strictly increasing chain
$(\F_\alpha)_{\alpha<\omega_1}$ of countably based Schur filters whose union
is not almost Schur.
\end{proposition}

\begin{proof}
Fix an admissible $u$.  The trace filter $\W_u$ has a base of cardinality at
most the continuum: the finite test lists range over $\ell_\infty$, and the
other parameters may be restricted to rational tolerances and integer tails.
It cannot have a countable base, because the source proves that every
countably based filter is strongly diagonal and hence Schur, whereas $\W_u$
is not even almost Schur \cite[Corollary 4.5 and Theorem 3.5]{source}.  Under
CH its character is therefore exactly $\omega_1$.

Enumerate a base as $(B_\xi)_{\xi<\omega_1}$.  Recursively let
$\F_{\alpha+1}$ be generated by $\F_\alpha$, $B_\alpha$, and one further
set $D_\alpha\in\W_u\setminus\F_\alpha$; at a countable limit take the union.
Such a $D_\alpha$ always exists, since the filter generated before stage
$\omega_1$ has a countable base and therefore cannot equal $\W_u$.  Every
$\F_\alpha$ is a proper subfilter of $\W_u$, has a countable base, and hence
is Schur by the cited source results.  The choices of $D_\alpha$ make the
chain strict, while the enumeration of the $B_\xi$ gives
\[
 \bigcup_{\alpha<\omega_1}\F_\alpha=\W_u,
\]
which is not almost Schur, witnessed by $(u_n)$.
\end{proof}

The proposition does not show that no basis for almost-Schur filters exists;
the source request permits many meanings of ``explicit.''  It shows sharply
why maximal almost-Schur filters cannot simply be obtained by the standard
chain-union form of Zorn's lemma.  Attempts through Schur dominators,
block-respecting filters, diagonal filters, and simultaneous blockers all
lead to the same unresolved extension issue.  Accordingly no answer to the
positive half is claimed.

\section*{Verification and literature boundary}

The proof uses only normalization, norm-density of $S_{\Q}$, and the source's
countable-base Schur theorem for the optional CH obstruction.  Each inclusion
direction in Theorem \ref{thm:basis} was checked directly from the definition
of filter convergence.  There is no computational dependency.

The bounded novelty audit searched the run indexes, the exact problem phrase,
the exact title and citations, and the phrases ``almost Schur filter'' and
``basis of filters'' through 11 August 2026.  Later filter-convergence papers
were found, but no answer to this basis request or this rational witness basis.
This is a bounded audit, not an exhaustive priority claim.

\begin{thebibliography}{9}
\bibitem{source}
A. Avilés López, B. Cascales Salinas, V. Kadets, and A. Leonov,
\emph{The Schur $\ell_1$ theorem for filters},
Zh. Mat. Fiz. Anal. Geom. \textbf{3} (2007), 383--398;
arXiv:0903.0659.  Included as \texttt{source\_paper.pdf}.
\end{thebibliography}

\clearpage
\section*{Exact source excerpt: Definition 6.7 and the open problem (page 15)}
\begin{center}
\includegraphics[width=0.97\linewidth,keepaspectratio]{figures/open_problem_crop.png}
\end{center}

\end{document}
